\documentclass[11pt,a4paper]{article}
\usepackage{fontspec}
\usepackage{polyglossia}\setdefaultlanguage{french}
\setmainfont{TeX Gyre Heros}
\newfontfamily\symf{DejaVu Sans}
\newcommand{\sym}[1]{{\symf #1}}
\setmonofont{DejaVu Sans Mono}[Scale=0.85]
\usepackage[a4paper,margin=2.2cm,headheight=14pt]{geometry}
\usepackage{xcolor}
\definecolor{vert}{HTML}{2E7D32}
\definecolor{vertclair}{HTML}{E8F5E9}
\definecolor{bleu}{HTML}{1565C0}
\definecolor{bleuclair}{HTML}{E3F2FD}
\definecolor{orange}{HTML}{E65100}
\definecolor{orangeclair}{HTML}{FFF3E0}
\definecolor{rouge}{HTML}{C62828}
\definecolor{rougeclair}{HTML}{FFEBEE}
\definecolor{gris}{HTML}{555555}
\definecolor{grisclair}{HTML}{F4F4F4}
\usepackage{tcolorbox}\tcbuselibrary{skins,breakable}
\newtcolorbox{idee}[1][]{enhanced,breakable,colback=vertclair,colframe=vert,boxrule=0.6pt,arc=2pt,left=8pt,right=8pt,top=6pt,bottom=6pt,fonttitle=\bfseries,title={#1}}
\newtcolorbox{exemple}[1][Exemple]{enhanced,breakable,colback=bleuclair,colframe=bleu,boxrule=0.6pt,arc=2pt,left=8pt,right=8pt,top=6pt,bottom=6pt,fonttitle=\bfseries,title={#1}}
\newtcolorbox{attention}[1][À savoir]{enhanced,breakable,colback=orangeclair,colframe=orange,boxrule=0.6pt,arc=2pt,left=8pt,right=8pt,top=6pt,bottom=6pt,fonttitle=\bfseries,title={#1}}
\newtcolorbox{probleme}[1][]{enhanced,breakable,colback=rougeclair,colframe=rouge,boxrule=0.6pt,arc=2pt,left=8pt,right=8pt,top=6pt,bottom=6pt,fonttitle=\bfseries,title={#1}}
\newtcolorbox{commande}{enhanced,breakable,colback=grisclair,colframe=gris!50,boxrule=0.4pt,arc=2pt,left=6pt,right=6pt,top=4pt,bottom=4pt,fontupper=\ttfamily\small}
\usepackage{tikz}\usetikzlibrary{positioning,arrows.meta,shapes.geometric,calc}
\usepackage{booktabs,tabularx,longtable,array,colortbl}
\renewcommand{\arraystretch}{1.25}
\newcolumntype{Y}{>{\raggedright\arraybackslash}X}
\newcolumntype{L}[1]{>{\raggedright\arraybackslash}p{#1}}
\usepackage{enumitem}\setlist{itemsep=2pt,topsep=4pt}
\usepackage{amsmath}
\usepackage{titlesec}
\titleformat{\section}{\Large\bfseries\color{vert}}{\thesection.}{0.6em}{}
\titleformat{\subsection}{\large\bfseries\color{vert!80!black}}{\thesubsection}{0.6em}{}
\titlespacing*{\section}{0pt}{16pt}{8pt}
\usepackage{fancyhdr}
\pagestyle{fancy}\fancyhf{}
\renewcommand{\headrulewidth}{0.3pt}
\usepackage{lastpage}
\cfoot{\normalfont\small\color{gris}\thepage\ / \pageref{LastPage}}
\usepackage[hidelinks]{hyperref}
\setlength{\parindent}{0pt}\setlength{\parskip}{5pt}
\usepackage{fvextra}
\DefineVerbatimEnvironment{code}{Verbatim}{fontsize=\small,breaklines,breakanywhere,frame=single,rulecolor=\color{gris!40},framesep=6pt,xleftmargin=2pt,xrightmargin=2pt,commandchars=\\\{\}}

\lhead{\normalfont\small\color{gris}AgroRot — Comment l'outil propose une rotation}
\rhead{\normalfont\small\color{gris}NASA Space Apps 2026 · Field Shift}

\begin{document}

\begin{titlepage}
\vspace*{2.5cm}
{\color{vert}\Huge\bfseries AgroRot}\par\vspace{6pt}
{\LARGE Comment l'outil propose une rotation de cultures}\par\vspace{14pt}
{\large Les données utilisées, les règles, les notes et les explications,\\ expliquées pour tout le monde}\par
\vspace{1.5cm}
\begin{idee}[Ce document en une phrase]
AgroRot compare la rotation actuelle d'un agriculteur avec des milliers d'autres possibles, élimine celles qui
enfreignent une règle agronomique, note les autres sur six critères à partir de 30 ans de données NASA et de ce
que l'agriculteur a saisi, puis présente les deux meilleures en expliquant pourquoi, avec des chiffres.
\end{idee}
\vfill
{\color{gris}
Projet : NASA Space Apps Challenge 2026, défi « Field Shift »\\
Région pilote : Yaoundé (Cameroun) — parcelle de démonstration à Mbankomo\\
Version du moteur : novembre 2026}
\end{titlepage}

\tableofcontents
\newpage

% =====================================================================
\section{L'idée en une page}

Une \textbf{rotation}, c'est l'ordre dans lequel on fait se succéder les cultures sur une même parcelle, saison
après saison. Une bonne rotation garde le sol fertile, casse le cycle des maladies, protège le sol contre
l'érosion et réduit le risque de manquer d'eau. Une mauvaise rotation (par exemple du maïs saison après saison)
épuise le sol et nourrit les mêmes parasites.

AgroRot aide l'agriculteur à choisir. Pour une parcelle donnée, il affiche \textbf{trois rotations côte à côte} :

\begin{center}
\begin{tabularx}{0.95\textwidth}{L{3.3cm}Y}
\toprule
\textbf{Rotation} & \textbf{Ce que c'est} \\
\midrule
Actuelle (référence) & Ce que l'agriculteur fait aujourd'hui (son historique répété, ou la rotation qu'il prévoit). Elle sert de point de comparaison. \\
Recommandée & La meilleure rotation trouvée, compte tenu des priorités de l'agriculteur. \\
Alternative & Une autre très bonne rotation, \emph{vraiment différente} de la recommandée, pour laisser un vrai choix. \\
\bottomrule
\end{tabularx}
\end{center}

Chaque rotation reçoit une \textbf{note sur 100}, le détail de six notes (climat, eau, fertilité, érosion,
maladies, priorités), un \textbf{indice de résilience} (comment elle tient si le climat se dégrade), un
\textbf{niveau de confiance} et des \textbf{explications} en phrases simples.

\begin{attention}[Ce que l'outil n'est pas]
\begin{itemize}
\item Ce n'est \textbf{pas une « intelligence artificielle » boîte noire} : il n'y a pas d'apprentissage automatique.
C'est un moteur de \textbf{règles agronomiques} publiées (FAO, NASA) ; chaque note se recalcule à la main et chaque
phrase d'explication vient d'une règle et d'un chiffre.
\item Il \textbf{ne remplace pas un agronome} local : il propose, l'agriculteur et son conseiller décident.
\item Il \textbf{ne prédit pas la météo de l'an prochain} : il juge une rotation sur le climat des 30 dernières
années (ce qui arrive « en général » et « les mauvaises années »), puis la teste sous des climats plus durs.
\end{itemize}
\end{attention}

% =====================================================================
\section{Petit vocabulaire}

\begin{center}
\begin{tabularx}{\textwidth}{L{3.6cm}Y}
\toprule
\textbf{Mot} & \textbf{Sens dans AgroRot} \\
\midrule
Saison S1 / S2 & À Yaoundé, deux saisons des pluies par an : \textbf{S1} (semis en mars, pluies jusqu'en juin) et \textbf{S2} (semis en août, pluies jusqu'en novembre). Une rotation sur 6 saisons couvre donc 3 ans. \\
Culture longue & Culture de 200 jours ou plus (manioc, igname…). Elle occupe plusieurs saisons d'affilée ; on l'affiche « suite » sur la saison suivante. \\
Jachère & Saison de repos, sans culture. \\
Plante de service & Plante semée pour le sol et non pour la récolte : mucuna, crotalaire (couvrent le sol, apportent de l'azote, réduisent certains parasites). \\
Légumineuse & Famille de plantes (arachide, niébé, haricot, soja, mucuna…) qui fixent l'azote de l'air et enrichissent le sol. \\
Famille botanique & Les plantes d'une même famille partagent souvent les mêmes maladies : on évite de les enchaîner. \\
ET0 & « Évapotranspiration de référence » : l'eau que l'air « pompe » chaque jour au sol et aux plantes (mm/jour). Plus il fait chaud, sec et venteux, plus elle est forte. \\
Kc & Coefficient de la culture : une culture en pleine croissance consomme \(K_c \times \text{ET0}\). Le Kc est faible au semis, fort à la floraison. \\
Bioagresseur & Maladie, ravageur, nématode ou plante parasite qui attaque les cultures. \\
\bottomrule
\end{tabularx}
\end{center}

% =====================================================================
\section{Les données utilisées}

AgroRot croise trois types de données : \textbf{le climat et le terrain} (satellites et modèles NASA),
\textbf{les connaissances sur les cultures} (FAO et littérature agronomique), et \textbf{ce que l'agriculteur sait de sa parcelle}.

\subsection{Climat et terrain (NASA)}

\begin{center}
\begin{tabularx}{\textwidth}{L{3.1cm}L{4.4cm}Y}
\toprule
\textbf{Source} & \textbf{Ce qu'on prend} & \textbf{À quoi ça sert} \\
\midrule
\textbf{NASA POWER} (point journalier, communauté « AG ») & Température moyenne, minimale et maximale, pluie, humidité, vent, rayonnement solaire, \textbf{chaque jour de 1996 à 2025} (30 ans, environ 11\,000 jours) au point de la parcelle & Base de tout le calcul : climat mois par mois, années sèches, jours de forte chaleur, pluies intenses, ET0, détection des saisons \\
ET0 Penman-Monteith (méthode FAO-56) & Calculée par AgroRot à partir des données POWER (méthode simplifiée de Hargreaves si un paramètre manque) & Besoin en eau des cultures (critère Eau) \\
\textbf{SRTM} (relief) & Altitude et pente de la parcelle & Altitude maximale de chaque culture ; risque d'érosion (critère Érosion) \\
\textbf{SMAP} (humidité du sol) et \textbf{MODIS NDVI} (vigueur de la végétation) & Valeurs mensuelles pré-téléchargées sur la zone pilote & Affichées dans le diagnostic pour aider l'agriculteur à comparer avec ce qu'il voit ; elles ne changent pas les notes \\
Prévisions météo (Meteomatics, sinon Open-Meteo) & 30 derniers jours + 10 prochains jours & Uniquement pour la \textbf{surveillance} après le choix d'une rotation (partie~\ref{sec:surveillance}) \\
\bottomrule
\end{tabularx}
\end{center}

\begin{attention}[Si Internet ou NASA POWER ne répondent pas]
Le service de données applique un « plan B » : il utilise d'abord la copie en cache (moins de 7 jours), sinon une
copie plus ancienne, sinon une \textbf{vraie série NASA POWER de Yaoundé (1996–2025) livrée avec l'application}.
L'écran affiche alors un bandeau (« données en cache » / « données locales ») et la confiance baisse un peu.
Le calcul fonctionne toujours.
\end{attention}

\subsection{Connaissances sur les cultures}

\begin{center}
\begin{tabularx}{\textwidth}{L{3.8cm}Y}
\toprule
\textbf{Fichier} & \textbf{Contenu et sources} \\
\midrule
Fiche culture (\texttt{crops\_cm.csv}) & \textbf{28 cultures} du Cameroun (maïs, manioc, arachide, niébé, haricot, igname, macabo, patate douce, plantain, tomate, piment, gombo, folon, ndolè, soja, sésame, riz, coton, cacao, palmier…, plus mucuna et crotalaire). Pour chacune : températures et pluies optimales et limites (\textbf{FAO ECOCROP}), durée du cycle, stades et Kc (\textbf{FAO-56}), pH, drainage, profondeur des racines, besoin en azote, azote fixé, résidus, couverture du sol, sensibilité à la sécheresse, délai de retour, usages (vente, alimentation, fourrage), bioagresseurs principaux. Une colonne « à vérifier » signale les valeurs estimées. 16 de ces cultures entrent en rotation à Yaoundé (les arbres comme le cacao ou le palmier n'entrent pas en rotation). \\
Table des bioagresseurs & Pour chaque maladie ou parasite noté : les cultures hôtes (fortes ou modérées), la gravité, la persistance dans le sol ou les résidus, et l'effet des cultures « assainissantes » (ex. : la crotalaire réduit les nématodes à galles). \\
Calendrier régional & Mois de semis et de fin des pluies de S1 et S2 à Yaoundé (FAO Crop Calendar, à confirmer avec l'IRAD). Les saisons détectées dans les données NASA ne l'ajustent que si elles restent cohérentes (± 1 mois). \\
Paramètres du moteur (\texttt{parametres.yaml}) & Tous les seuils et coefficients des règles, réunis dans un seul fichier lisible : aucune valeur agronomique n'est cachée dans le code. \\
\bottomrule
\end{tabularx}
\end{center}

\subsection{Ce que l'agriculteur saisit}

\begin{center}
\begin{tabularx}{\textwidth}{L{3.4cm}Y}
\toprule
\textbf{Saisie} & \textbf{Utilisation} \\
\midrule
Contour de la parcelle (points sur la carte) & Surface, centre (pour aller chercher le climat NASA de ce point exact), altitude et pente \\
Sol : texture, pH, drainage, matière organique, pente & Règles de sol (pH, drainage), réserve d'eau du sol (selon la texture), érosion ; chaque valeur manquante fait baisser la confiance \\
Eau : irrigation, eau disponible en saison sèche & Bilan d'eau (critère Eau) \\
Pratiques : engrais, part des résidus laissés au champ & Fertilité (critère Sol) \\
Historique : cultures des 3 dernières saisons & Construit la rotation « actuelle » et sert de point de départ pour les délais de retour et les maladies \\
Maladies déjà observées & Pression de départ des bioagresseurs \\
Priorités (classement de 6 objectifs) & Poids des critères (partie~\ref{sec:note}) \\
Cultures exclues, cultures qui se vendent localement & Une culture exclue n'est jamais proposée ; le revenu ne compte que pour les cultures qui ont un débouché \\
Corrections du climat (« il pleut moins en juillet chez moi ») & Ajustent la série NASA mois par mois (pluie en \%, température en °C) \\
\bottomrule
\end{tabularx}
\end{center}

% =====================================================================
\clearpage
\section{Le chemin d'un calcul}

Quand l'agriculteur clique sur « Calculer », voici ce qui se passe, en un peu plus d'une seconde :

\begin{center}
\begin{tikzpicture}[
  box/.style={rectangle,rounded corners=3pt,draw=vert,fill=vertclair,thick,minimum width=4.3cm,minimum height=1cm,align=center,font=\small},
  data/.style={rectangle,rounded corners=3pt,draw=bleu,fill=bleuclair,thick,minimum width=3.6cm,minimum height=0.9cm,align=center,font=\small},
  fl/.style={-{Latex[length=2.5mm]},thick,gris}]
\node[data] (n) {Climat NASA\\(30 ans, au point)};
\node[data,right=0.6cm of n] (f) {Fiche culture,\\bioagresseurs};
\node[data,right=0.6cm of f] (a) {Saisies de\\l'agriculteur};
\node[box,below=0.9cm of f] (e1) {\textbf{1. Filtre d'aptitude}\\cultures possibles ici};
\node[box,below=0.5cm of e1] (e2) {\textbf{2. Générateur}\\des milliers de rotations};
\node[box,below=0.5cm of e2] (e3) {\textbf{3. Règles éliminatoires}\\C0 à C7};
\node[box,below=0.5cm of e3] (e4) {\textbf{4. Six notes} S1 à S6\\+ note globale sur 100};
\node[box,below=0.5cm of e4] (e5) {\textbf{5. Choix} : recommandée\\+ alternative différente};
\node[box,below=0.5cm of e5] (e6) {\textbf{6. Résilience, confiance,}\\\textbf{explications}};
\draw[fl] (n.south) -- (e1.north west);
\draw[fl] (f.south) -- (e1.north);
\draw[fl] (a.south) -- (e1.north east);
\draw[fl] (e1) -- (e2); \draw[fl] (e2) -- (e3); \draw[fl] (e3) -- (e4); \draw[fl] (e4) -- (e5); \draw[fl] (e5) -- (e6);
\draw[fl,dashed] (e4.east) -- ++(1.5,0) |- node[pos=0.25,right,font=\footnotesize,align=left]{boucle :\\on garde\\les meilleures} (e2.east);
\end{tikzpicture}
\end{center}

La rotation actuelle de l'agriculteur passe par les mêmes étapes 3 à 6, pour être comparée équitablement.

% =====================================================================
\section{Étape 1 — Quelles cultures sont possibles ici ?}

Pour chaque culture, on compare le climat moyen de la parcelle (température annuelle, pluie annuelle) à ce que la
culture aime, d'après ECOCROP. On calcule une \textbf{aptitude} entre 0 et 1 :

\begin{itemize}
\item 1 si la température et la pluie sont dans la plage idéale de la culture ;
\item qui diminue progressivement quand on s'approche des limites ;
\item 0 au-delà des limites, ou si la parcelle est plus haute que l'altitude maximale de la culture.
\end{itemize}

Les cultures dont l'aptitude est \textbf{inférieure à 0,3} sont écartées dès le départ. Avec une irrigation
déclarée, la pluie n'est plus un facteur limitant (aptitude « pluie » d'au moins 0,8). Pour les cultures très
courtes (60 jours ou moins, comme le folon), seule la limite basse de pluie compte : elles sont cultivées dans une
seule saison et les limites « annuelles » d'ECOCROP n'ont pas de sens pour elles.

% =====================================================================
\section{Étape 2 — Fabriquer des rotations à comparer}

Le nombre de rotations possibles est énorme : avec 10 cultures candidates par saison sur 6 saisons, cela fait un
million de combinaisons. Le moteur ne les essaie pas toutes : il construit les rotations \textbf{saison par
saison} et, à chaque saison, ne garde que les \textbf{120 meilleurs débuts de rotation} (technique dite « recherche
en faisceau »). C'est comme un concours avec des éliminatoires à chaque tour.

Deux garde-fous évitent des propositions absurdes :
\begin{itemize}
\item au plus \textbf{une saison sans récolte} (jachère ou plante de service) toutes les 4 saisons : l'agriculteur doit pouvoir manger et vendre ;
\item une culture longue (manioc) occupe automatiquement les saisons qu'elle couvre.
\end{itemize}

% =====================================================================
\section{Étape 3 — Les règles éliminatoires}

Certaines erreurs sont rédhibitoires. Une rotation qui enfreint une \textbf{règle dure} n'est jamais proposée.
La rotation actuelle, elle, est notée même si elle enfreint des règles, mais elle perd \textbf{10 \% de sa note par
règle dure enfreinte} (une fois par culture concernée, au maximum 50 \%), et l'explication le dit.

\begin{center}
\begin{tabularx}{\textwidth}{L{1cm}L{3.4cm}YL{1.6cm}}
\toprule
\textbf{Code} & \textbf{Règle} & \textbf{Ce qu'on vérifie} & \textbf{Type} \\
\midrule
C0 & Culture exclue & L'agriculteur a dit qu'il ne veut pas (ou ne peut pas) faire cette culture & Absolue \\
C1 & Climat & Nuits trop froides ou jours trop chauds pour la culture (on regarde les 10 \% d'années les plus froides et les 10 \% les plus chaudes), pluie annuelle hors limites & Dure \\
C2 & Saison & La culture doit tenir dans la saison des pluies où elle est semée. Exception : une culture longue peu sensible à la sécheresse (manioc) peut traverser la saison sèche & Dure \\
C3 & Chevauchement & On ne peut pas semer sur une parcelle encore occupée par une culture longue & Dure \\
C4 & Délai de retour & Une culture ne revient pas sur la même parcelle avant son délai (ex. : maïs 1 an, manioc 2 ans) & Dure \\
C5 & Même famille & Deux cultures de la même famille à la suite (maïs puis maïs, ou maïs puis riz) & Souple : pénalisée dans la note Maladies \\
C6 & Sol & pH hors de la plage de la culture (avec une tolérance de 0,5), ou sol mal drainé pour une culture qui ne le supporte pas & Dure \\
C7 & Pression de parasites & Un bioagresseur atteint un niveau critique au fil des saisons & Souple : alerte \\
\bottomrule
\end{tabularx}
\end{center}

Si aucune rotation ne respecte toutes les règles (par exemple quand l'agriculteur a exclu presque toutes les
cultures), le moteur relâche les règles C4 et C5 et l'affiche clairement (« règles secondaires assouplies »).

% =====================================================================
\section{Étape 4 — Les six notes}

Chaque rotation qui passe les règles reçoit six notes sur 100. Chacune répond à une question simple.

\subsection*{S1 — Climat : « Les cultures sont-elles dans leurs bonnes températures ? »}
Pour chaque culture, on mesure la part de son cycle passée dans sa plage de température idéale (ECOCROP), avec une
baisse progressive vers les limites. On retire une pénalité pour les \textbf{jours de forte chaleur} pendant la
floraison (plus de 5 °C au-dessus de l'optimum), comptés sur les 30 années.

\subsection*{S2 — Eau : « Combien d'années sur 10 la pluie couvre-t-elle le besoin des cultures ? »}
C'est le cœur du calcul climatique, selon la méthode FAO-56 :
\begin{enumerate}
\item besoin de la culture, mois par mois : \(K_c \times \text{ET0}\) (le Kc suit les stades : semis, croissance, floraison, maturité) ;
\item eau disponible : 80 \% de la pluie (le reste ruisselle ou s'évapore) + l'irrigation éventuelle ;
\item \textbf{réserve du sol} : l'eau en trop d'un mois humide reste disponible les mois suivants, dans la limite de
ce que le sol peut retenir (selon la texture) dans la zone des racines ;
\item on refait ce calcul pour \textbf{chacune des 30 années} et on compte les années « réussies », où au moins
80 \% du besoin est couvert (65 \% pour les cultures peu sensibles à la sécheresse comme le manioc ou la patate
douce, qui réduisent leur feuillage puis repartent).
\end{enumerate}
La note est d'autant plus basse que les années ratées sont fréquentes et que la culture est sensible à la sécheresse.

\begin{exemple}[Exemple]
Le maïs semé en S1 à Mbankomo couvre son besoin 29 années sur 30 : excellent. Le macabo semé en S2 traverse la
saison sèche de décembre à février et ne réussit presque jamais sans irrigation : mauvaise note Eau, et en plus la
règle C2 l'élimine.
\end{exemple}

\subsection*{S3 — Fertilité du sol : « La rotation enrichit-elle ou épuise-t-elle le sol ? »}
Un bilan par points, saison par saison :
\begin{itemize}
\item \textbf{azote} : +3 pour une légumineuse ; −1, −2 ou −3 selon le besoin en azote de la culture ; un bonus
quand une légumineuse précède une culture exigeante (arachide puis maïs) ; une pénalité quand deux cultures
exigeantes se suivent (moins forte si l'agriculteur met de l'engrais) ;
\item \textbf{matière organique} : points selon la quantité de résidus de la culture et la part laissée au champ, plus un bonus pour la fumure organique.
\end{itemize}
Le total (60 \% azote, 40 \% matière organique) est ramené sur 100.

\subsection*{S4 — Érosion : « Le sol est-il couvert quand tombent les grosses pluies ? »}
Pour chaque mois de la rotation, on regarde la \textbf{part des pluies intenses de l'année} qui tombe ce mois-là
(calculée sur 30 ans de NASA POWER) et \textbf{à quel point le sol est couvert} (sol nu, culture jeune, culture
développée, plante de couverture). Un mois de grosses pluies sur sol nu coûte cher, surtout si la pente est forte
(version simplifiée de l'équation universelle de perte en terre). Les explications indiquent le mois le plus risqué
et suggèrent une plante de couverture.

\subsection*{S5 — Maladies et ravageurs : « Les parasites ont-ils de quoi se nourrir saison après saison ? »}
Pour chaque bioagresseur, on simule une \textbf{pression} qui :
\begin{itemize}
\item \textbf{monte} quand une culture hôte est en place (fortement pour un hôte fort, moins pour un hôte modéré) ;
\item \textbf{redescend} doucement quand il n'y a pas d'hôte (vite pour un parasite de courte durée, très lentement pour un parasite qui survit des années dans le sol) ;
\item \textbf{chute} quand une culture assainissante passe (crotalaire contre les nématodes, par exemple) ;
\item part d'un niveau plus élevé si l'agriculteur a déjà observé ce problème.
\end{itemize}
On simule la rotation deux fois de suite et on note le deuxième passage (le régime « installé »). Plus la pression
totale, pondérée par la gravité, est forte, plus la note baisse. Enchaîner deux cultures de la même famille (C5)
ajoute une pénalité.

\subsection*{S6 — Vos priorités : « La rotation sert-elle ce que l'agriculteur veut ? »}
Part des saisons qui portent une culture \textbf{vendable sur ses marchés} (si le revenu compte pour lui), une
culture qui donne du \textbf{fourrage} (si l'élevage compte), et un peu de poids pour l'alimentation familiale.
Une culture « vendable » sans débouché local déclaré ne compte pas pour le revenu, et l'explication le signale.

% =====================================================================
\section{La note globale et le rôle des priorités}\label{sec:note}

La note globale est une \textbf{moyenne pondérée} des six notes :
\[
G = w_1 S_1 + w_2 S_2 + w_3 S_3 + w_4 S_4 + w_5 S_5 + w_6 S_6 \qquad (w_1 + \dots + w_6 = 1)
\]
puis la pénalité éventuelle des règles dures (étape 3).

Les poids viennent des \textbf{priorités} de l'agriculteur. Il classe six objectifs (économiser l'eau, revenu,
santé du sol, azote, lutte contre l'érosion, fourrage) ; le classement est transformé en poids par une méthode
standard d'aide à la décision (« Rank Order Centroid ») : le premier objectif pèse 41 \%, le deuxième 24 \%, puis
16 \%, 10 \%, 6 \% et 3 \%. Ces objectifs sont ensuite répartis sur les critères :

\begin{center}
\begin{tabular}{ll}
\toprule
\textbf{Objectif de l'agriculteur} & \textbf{Critère qui en profite} \\
\midrule
Économiser l'eau & S2 Eau \\
Revenu, fourrage & S6 Priorités \\
Azote & S3 Fertilité \\
Santé du sol & moitié S3 Fertilité, moitié S5 Maladies \\
Lutte contre l'érosion & S4 Érosion \\
\bottomrule
\end{tabular}
\end{center}

\begin{attention}[Deux critères protégés]
Le climat (S1) et les maladies (S5) gardent toujours \textbf{au moins 15 \%} chacun, quelles que soient les
priorités : une rotation qui ignore le climat ou qui nourrit les parasites ne doit jamais sortir en tête parce que
l'agriculteur a tout misé sur le revenu. Sans priorités saisies, les six objectifs pèsent autant.
\end{attention}

\begin{exemple}[Exemple : l'agriculteur classe « économiser l'eau » en premier, puis revenu, santé du sol, azote, érosion, fourrage]
Poids obtenus : Eau 29 \%, Maladies 20 \%, Priorités 19 \%, Climat 15 \%, Fertilité 13 \%, Érosion 4 \%.
\end{exemple}

% =====================================================================
\section{Étape 5 — Choisir la recommandée et l'alternative}

\begin{itemize}
\item La \textbf{recommandée} est la rotation qui a la meilleure note globale parmi celles qui respectent toutes les règles dures.
\item L'\textbf{alternative} est la meilleure rotation qui diffère de la recommandée sur \textbf{au moins la moitié
des saisons} et qui utilise \textbf{au moins un tiers de cultures différentes}. Sans cette condition, l'alternative
serait souvent la même rotation décalée d'une saison, ce qui n'offre aucun vrai choix.
\item Chacune reçoit un nom qui dit son point fort (« économie d'eau d'abord », « santé des cultures »…).
\end{itemize}

L'agriculteur peut ensuite \textbf{remplacer une culture à la main} : la rotation modifiée est recalculée, les
règles enfreintes sont signalées et la note est mise à jour. Il garde la main.

% =====================================================================
\section{Étape 6 — Résilience, confiance et explications}

\subsection{Résilience : « Et si le climat se dégrade ? »}
On recalcule la rotation sous trois climats plus durs : \textbf{20 \% de pluie en moins}, \textbf{2 °C de plus},
et \textbf{les deux à la fois}. L'indice de résilience est la note dans le pire de ces trois cas divisée par la
note normale (plafonné à 1). Un indice de 0,9 signifie que la rotation garde 90 \% de sa note dans le pire cas.

Indépendamment, l'écran « Et si… » permet à l'agriculteur de tester son propre scénario : pluie de −30 \% à
+30 \%, température de 0 à +3 °C, horizon 0, 5 ou 10 ans. La chaleur augmente aussi l'ET0 (environ +3 \% par °C).

\subsection{Confiance}
La confiance (élevée, moyenne, faible) dit \textbf{à quel point on peut se fier au résultat}, pas s'il est bon.
Elle part de la qualité des informations sur le sol (valeur mesurée, estimée ou absente) et baisse quand :
les données climatiques viennent du secours plutôt que de NASA POWER en direct, des valeurs de la fiche culture
sont marquées « à vérifier », ou les pratiques ne sont pas renseignées. Élevée à partir de 0,80, moyenne à partir
de 0,60.

\subsection{Explications}
Chaque calcul garde la trace des règles qui ont pesé, avec leurs chiffres. Le moteur les trie par importance et
les range en cinq rubriques :

\begin{center}
\begin{tabularx}{\textwidth}{L{3.6cm}Y}
\toprule
\textbf{Rubrique} & \textbf{Exemple réel (Mbankomo)} \\
\midrule
Pourquoi & « Fertilité du sol (S3) : 48 contre 26 pour votre rotation actuelle. » \\
Preuves & « Bilan hydrique FAO-56 (Kc × ET0, 80 \% de la pluie efficace) sur NASA POWER, 1996–2025. » \\
Bénéfices & « Arachide (légumineuse) avant Maïs : apport d'azote pour une culture exigeante (environ 40 à 100 kg N/ha selon la littérature). » \\
Informations manquantes & « Présence réelle de bioagresseurs non renseignée : pression initiale supposée nulle. » \\
Risques à vérifier sur place & « Sol nu en juillet, mois qui reçoit 7 \% des pluies intenses de l'année (pente 6 \%). Piste : une couverture (Mucuna) à ce moment-là. » \\
\bottomrule
\end{tabularx}
\end{center}

% =====================================================================
\section{Un exemple complet : la parcelle de Mbankomo}

Parcelle de démonstration près de Yaoundé : sol argilo-limoneux, pH 5,8, pente d'environ 3°, pas d'irrigation,
historique « maïs, maïs, manioc », priorité « économiser l'eau », gombo exclu. Climat : vraie série NASA POWER
1996–2025 (environ 1\,460 mm de pluie par an). Rotation sur 6 saisons (3 ans).

\begin{center}
\small
\begin{tabularx}{\textwidth}{L{2.6cm}Ycccccccc}
\toprule
\textbf{Rotation} & \textbf{Saisons 1 à 6} & \textbf{S1} & \textbf{S2} & \textbf{S3} & \textbf{S4} & \textbf{S5} & \textbf{S6} & \textbf{Note} & \textbf{Résil.} \\
\midrule
Actuelle & maïs, maïs, manioc, (suite), maïs, manioc & 100 & 97 & 26 & 83 & 64 & 96 & \textbf{65} & 0,84 \\
Recommandée & arachide, maïs, mucuna, maïs, manioc, (suite) & 100 & 99 & 48 & 85 & 81 & 82 & \textbf{85} & 0,88 \\
Alternative & maïs, patate douce, crotalaire, maïs, manioc, (suite) & 100 & 99 & 41 & 84 & 84 & 81 & \textbf{85} & 0,89 \\
\bottomrule
\end{tabularx}
\end{center}

\textbf{Comment lire ce tableau :}
\begin{itemize}
\item La rotation actuelle est très bonne pour le revenu (S6 = 96 : maïs et manioc se vendent) mais \textbf{épuise
le sol} (S3 = 26) et \textbf{nourrit les parasites} (S5 = 64). Elle enfreint aussi des délais de retour (le maïs
revient après 5 mois au lieu d'un an ; le manioc revient trop vite) : 2 cultures concernées, donc −20 \% sur sa note.
\item La recommandée place des \textbf{légumineuses avant le maïs} (arachide, mucuna) : la fertilité double presque
et les maladies reculent. Elle perd un peu sur le revenu (une saison de mucuna ne se vend pas) mais gagne 20 points.
\item L'alternative fait un autre choix (patate douce, crotalaire contre les nématodes) pour une note équivalente.
\end{itemize}

\begin{exemple}[Et si la pluie baisse de 30 \% et la température monte de 2 °C ?]
La rotation actuelle tombe à \textbf{50} (sa note Eau passe de 97 à 67). Le moteur propose alors une
recommandée adaptée (niébé, maïs, mucuna, maïs, niébé, maïs) qui garde \textbf{85}. C'est le message du défi
« Field Shift » : changer de rotation permet de s'adapter à un climat qui change.
\end{exemple}

% =====================================================================
\section{Après le choix : la surveillance de la rotation}\label{sec:surveillance}

Quand l'agriculteur sauvegarde une rotation, elle devient la rotation \textbf{surveillée} de la parcelle. Chaque
jour, l'outil compare la \textbf{météo des 30 derniers jours et la prévision des 10 prochains jours} à ce que la
rotation suppose. Il sait où on en est (avant le semis, fenêtre de semis, culture en place, entre deux cultures)
et envoie une notification à tous les membres de la ferme (agriculteur, co-exploitant, conseiller) en cas de besoin :

\begin{center}
\begin{tabularx}{\textwidth}{L{4.2cm}Y}
\toprule
\textbf{Notification} & \textbf{Quand} \\
\midrule
Bonnes conditions pour semer / Retarder le semis & Autour de la date de semis prévue : au moins 40 mm de pluie (10 jours passés + 7 jours prévus), ou moins de 20 mm \\
Manque d'eau probable & Culture en place sans irrigation : sur 3 semaines, le besoin de la culture dépasse nettement la pluie (seuil plus bas pour les cultures sensibles à la sécheresse) \\
Forte pluie prévue & Un jour prévu à 50 mm ou plus (critique à 80 mm), avec des conseils selon la pente et la culture \\
Fortes chaleurs & Au moins 2 jours prévus au-dessus de ce que supporte la culture \\
Saison anormale / Rotation plus adaptée & Pluie du dernier mois sous 60 \% ou au-dessus de 160 \% de la normale : le moteur recalcule la rotation « si la saison continue ainsi ». Si la note baisse de 10 points ou plus, il propose de relancer le calcul. \\
\bottomrule
\end{tabularx}
\end{center}

Une même alerte n'est pas renvoyée deux fois en 7 jours. Chaque notification contient des conseils concrets
(pailler, désherber, curer les rigoles, attendre une pluie de 20 mm avant de semer…).

% =====================================================================
\section{Limites et honnêteté}

\begin{itemize}
\item \textbf{Précision du climat} : NASA POWER décrit une maille d'environ 50 km. Le relief local, un bas-fond ou
une forêt voisine peuvent changer les choses : c'est pourquoi l'agriculteur peut corriger la pluie et la température mois par mois.
\item \textbf{Fiche culture} : les valeurs viennent de bases internationales (ECOCROP, FAO-56) ; certaines sont
estimées et marquées « à vérifier ». Elles doivent être validées avec des agronomes locaux (IRAD, par exemple).
\item \textbf{Ce qui n'est pas modélisé} : les prix du marché, la main-d'œuvre disponible au jour le jour, les
variétés locales, les associations de cultures dans la même saison.
\item \textbf{Le jugement reste humain} : l'outil propose, explique et alerte ; l'agriculteur et son conseiller décident.
\end{itemize}

\begin{idee}[Pour aller plus loin (équipe technique)]
Les formules exactes sont dans le document technique « Principe de scoring » (LaTeX). Tous les seuils sont dans
\texttt{services/agrorot-core/config/parametres.yaml}, et les choix de calibrage sont expliqués dans
\texttt{services/agrorot-core/README.md}. Les règles de surveillance sont en tête de
\texttt{services/backend/app/domain/surveillance.py}.
\end{idee}

\end{document}
